\section{Preliminaries}
\label{sec:preliminaries}

% \subsection{Forward and inverse kinematics}
% \sbmp algorithms operate both in \cspace ~and in the workspace $E$. The functions translating configurations in \cspace, that is, vectors of \dof values, and points in $E$ and vice versa are the forward and inverse kinematics functions. Below are the definitions of forward and inverse kinematics we use in this paper. 

% \emph{\textbf{Forward kinematics (\fk).}}
% The function $\fk:\cspace \longrightarrow E$ takes a configuration $c\in \cspace$ as input and returns the position of $r$'s end-effector in $E$ when it assumes configuration $c$. Note that this position is not a full 6\dof pose, and in the context of visibility-based tools means the apex of the \fov $\C$.

% \emph{\textbf{Inverse kinematics (\ik).}}
% The function $\ik : E \longrightarrow \cspace$ takes a point $p\in E$ as input and returns a configuration $c$ of $r$ such that $\fk(c)=p$.

\subsection{Problem definition}
We formally define the \tvmp problem below. While this is, to the best of our knowledge, the first time the problem has been explicitly characterized in literature, it is a fundamental challenge inherent to many real-world applications.

% \paragraph*{\textbf{\tvmp problem definition}}
Let $r$ be a robot equipped with a visibility-based tool characterized by a conic (or conic frustum) \fov beam $\mathcal{C}$, and let $\visatcfg{x}$ denote the conic \fov emitted from $r$ at configuration $x$. Given an environment (workspace) $E \subseteq \R^3$, a target set $P\subseteq E$, and an initial configuration $s$ of $r$, the goal is to compute a motion plan that transitions $r$ from $s$ to some terminal configuration $t$, such that $P\subseteq \visatcfg{t}$. An illustration is given in Figure \ref{fig:tvmp}b. A natural variant of this problem, termed $\alpha$-\tvmp, requires that at least an $\alpha$-fraction of $P$ is contained within $\visatcfg{t}$, for some $\alpha \in (0, 1)$. Furthermore, this definition may extend to an optimization problem not included in the scope of this work: finding a configuration $t$ that maximizes the visible fraction of $P$.



\subsection{Simplifying assumptions}
In this work, we adopt several simplifying assumptions detailed below along with their underlying justifications. 

% \emph{\textbf{\fov parameterization.}}
% We model the robot's \fov $\C$ as a conic volume with its apex centered on a functional face of the visibility device. The cone is aligned with the face's normal vector and is geometrically defined by an aperture angle $\gamma$ and a range parameter $h$. This choice reflects the standard conic or rectangular frustum-shaped \fov{}s in most directional visibility-based devices. 

\emph{\textbf{Target set.}}
While the target set $P$ can be any arbitrary set of points in $E$, our algorithms first compute an approximation of the minimum enclosing sphere $P_S$ of $P$, and solve the \tvmp problem for $P_S$ rather than $P$. The reasoning behind this assumption is based on several considerations. 
\begin{enumerate}
    \item The most likely use-case for this algorithm may involve a single object or a compact region in $\R^3$.
    \item The assumption remains valid in the case of multiple target objects not occluded by obstacles.
    \item The case where the target objects are occluded by obstacles or are placed in close proximity to obstacles, e.g., in different compartments in a shelf or on a table, can be handled by solving an $\alpha$-\tvmp problem.
\end{enumerate}
We further assume that the points within $P$ are non-occluding. Given that $P$ is modeled as a sphere, this implies that all viewpoints yielding the same visible surface area are functionally equivalent. Consequently, the planning task does not require a specific orientation relative to the target.

\emph{\textbf{Environment.}}
We assume that the environment can be sufficiently approximated by basic geometries such as Axis-Aligned Bounding Boxes (\aabb), spheres, \new{or triangular meshes}. This assumption is required to leverage GPU-accelerated operations for processing large batches of visibility queries efficiently. Such approximations of environmental and robotic geometry are standard practice in motion planning literature to ensure computational tractability \cite{tkk-mmvsbp-23}.

\subsection{Visibility integrity (\vi)}
To efficiently handle visibility queries, \visprm maintains a tree data structure based on the \vi measure \cite{lk-fmtvw-16, zhvml-c3dfrvuvi-19}. The \vi score of a point set is an approximation of the intersection-over-union of their visibility polyhedrons. Thus, \vi measures the degree to which a set of points shares a common visible workspace (see Figure \ref{fig:vi:example}). Given a finite point set $X\subseteq \R^d$, we denote the subset of $X$ visible from a point $p\in \R^d$ by $V_X(p)\subseteq X$. The \vi of a point set $P\subseteq \R^d$ is defined by
\begin{equation}
    VI(P) = \frac{|\Set{\bigcap V_X(p)}{p\in P}|}{|\Set{\bigcup V_X(p)}{p\in P}|},
\end{equation}
where operator $|\cdot |$ denotes the cardinality of a discrete set. %While the extension to more general settings is simple, it is not required for our purposes and thus omitted for the sake of brevity.

% \begin{figure}[t!]
%     %\vspace{-0.3cm}
%     \centering
%     \begin{subfigure}{\linewidth}
%     \centering
%     \includegraphics[width=\linewidth,page=14]{figs/vi_example.pdf}
%         \caption{}
%         \label{fig:vi:example:high}
%     \end{subfigure}
%     \hfil
%     \begin{subfigure}{\linewidth}
%     \includegraphics[width=\linewidth,page=15]{figs/vi_example.pdf}
%         \caption{}
%         \label{fig:vi:example:low}
%     \end{subfigure}
%     %\vspace{-0.2cm}
%     \caption{Visibility examples in a planar environment represented by a set of grid points $X$ (black). (a) The blue points exhibit high \vi, as they share nearly identical visibility of the points in $X$. (b) The red points have divergent visibility of points in $X$ yielding a low \vi.}
%     \label{fig:vi:example}
%     %\vspace{-0.7cm}
% \end{figure}


\begin{figure}[!t]
\centering
\subfloat[]{\includegraphics[width=0.4\linewidth,page=14]{figs/vi_example.pdf}%
\label{fig:vi:example:high}}
\hfil
\subfloat[]{\includegraphics[width=0.4\linewidth,page=15]{figs/vi_example.pdf}%
\label{fig:vi:example:low}}
\caption{Visibility examples in a planar environment represented by a set of grid points $X$ (black). (a) The blue points exhibit high \vi, as they share nearly identical visibility of the points in $X$. (b) The red points have divergent visibility of points in $X$ yielding a low \vi.}
\label{fig:vi:example}
\vspace{-0.5cm}
\end{figure}